%!TeX root = main.tex
%!encoding = UTF-8 Unicode

\section{Figures of Merit for Electromagnetic Emission Spectra}
A figure of merit (FOM) is a single quantitative value used to characterize the performance of a system, device, or process with respect to a defined objective. 
It condenses relevant performance into a value that can be compared across designs or operating conditions. 
In this thesis, FOMs are used to assess and compare the EMC-relevant performance of switching transients. 
They may be derived from the time-domain waveform, the frequency-domain spectrum, or both, and applied to comparisons across operating points, parameter variations, or changes over time. 
When such a sweep includes many conditions, direct comparison of all corresponding spectra can become challenging, and FOMs provide a compact basis for comparison. 
However, a single value cannot represent all information in a waveform or spectrum: each FOM comes at the cost of a loss of information.

The FOMs considered here cover time-domain, hybrid, and frequency-domain characteristics. 
The Heisenberg-Gabor Inequality combines temporal and spectral characteristics, while the frequency-domain FOMs describe broadband behavior or focus on selected frequencies, including resonances between the semiconductor output capacitance and the parasitic inductance of the commutation cell.

\subsubsection*{Rise time {\boldmath\mi{t}{10-90\%}}}
For transients with the same overall shape, rise time provides a useful basis for comparison because the second spectral corner shifts to lower frequencies as the transition duration increases, as described by \eref{eq:fc2_rise_time} and illustrated in \fref{fig:Diss_SpectrumCreation_firstOrder}. 
However, rise time alone does not characterize waveform shape and is therefore insufficient for reliable comparisons between different transient types, such as active-off and passive-off events, or between operating conditions when a parameter change alters the transient shape. 
Piecewise-linear and Gaussian-smoothed transitions can have substantially different spectral decay even at the same rise time; the contrast is illustrated in \fref{fig:Diss_SpectrumCreation_Heisenberg}, and reflected in \eref{eq:gaussian_spectrum}.



\subsubsection*{Maximum Slew Rate {$\boldmath\sddt{v}_\mathrm{max}$}}
The maximum voltage slew rate, $\sddt{v}_\mathrm{max}$, quantifies the steepest rate of voltage change during a switching transient. 
However, it does not map directly to a spectral behavior; its value is influenced by the voltage change, related to the DC-link voltage \mi{V}{DC}, and by the transition duration, related to the rise time \mi{t}{10-90\%}. 
Like the rise time, it is therefore most useful for comparing transients of similar shape.

Although slew rate is commonly used to quantify switching-related EM emissions, the author considers it unsuitable for this purpose. 
Instead, it should only be used when specifying driver-isolation requirements. 
A voltage transient across the isolation capacitance generates a displacement current, $i = C_\mathrm{iso}\sfrac{\mathrm{d}v}{\mathrm{d}t}$, which can interfere with the driver's normal operation.

\subsubsection*{Heisenberg-Gabor Inequality (HGI)}
As discussed in Section~\ref{sec:Fundamental_Generation}, a Gaussian-smoothed switching transient has been proposed as a theoretically EM-noise-optimal waveform. 
The Heisenberg-Gabor Inequality (HGI) uses this optimum as a reference for assessing switching transients based on their temporal and spectral dispersion~\cite{patin_power_2015,martinez-padron_definition_2022}. 
The Gaussian smoothing function is
\begin{equation}
	g(t) = \frac{1}{\sigma_t \sqrt{2\pi}}\cdot e^{-\dfrac{t^2}{2\sigma_t^2}},
\end{equation}
where the resulting semiconductor voltage follows from the convolution
\begin{equation}
	\mi{v}{DS,opt} = \mi{V}{DC}\cdot\mathrm{rect}(\frac{t}{\mathrm{d}T_\mathrm{s}}) * g(t)
\end{equation}
The resulting waveform for $\sigma_t = \SI{10}{ns}$ is shown in Fig. \ref{fig:IdealVsReal}.

The dispersions are calculated from the transient function $\lambda = \frac{\text{d}\mi{v}{DS}}{\text{d}t}$, the derivative of the voltage transient.

\begin{equation}
	\sigma_t = \sqrt{\frac{\int (t-\mi{t}{b})^2 \cdot |\lambda(t)|^2 \text{d}t}{\int |\lambda(t)|^2 \text{d}t}}\text{, with } \mi{t}{b} = \int t\cdot\lambda(t) \text{d}t
\end{equation}

\begin{equation}
\sigma_f = \sqrt{\frac{\int (\omega-\mi{\omega}{b})^2 \cdot |\Lambda(t)|^2 \text{d}\omega}{\int |\Lambda(t)|^2 \text{d}\omega}}\text{, with } \mi{\omega}{b} = \int \omega\cdot\Lambda(t) \text{d}\omega
\end{equation}

Here, $\Lambda(\omega)$ is the Fourier Transform of $\lambda(t)$.
Both can be used to calculate the FOM 
\begin{equation}
	HGI = \sigma_f \cdot \sigma_t \geq 0.5.
\end{equation}

The lower bound of 0.5 is attained by a Gaussian waveform. 
HGI therefore indicates the time--frequency spread product of the transient relative to this theoretical limit; a larger value corresponds to a larger product, not necessarily to greater high-frequency emission in every frequency band.

\subsubsection*{Band-Averaged Spectral Magnitude {\boldmath$\bar{V}_{\Delta f}$}}
To quantify broadband emission over a selected frequency band, the frequency-bin-width-weighted mean of the spectral magnitudes is defined as:

\begin{equation}
	\bar{v}_{10-150 \mathrm{MHz}} = \frac{1}{\Delta f_{tot}} \sum_{n = n_{10 \mathrm{MHz}}}^{n_{150 \mathrm{MHz}}} \Delta f_n \cdot a_n
\end{equation}

Here, $a_n$ is the magnitude of the Fourier coefficient in frequency bin $n$, $\Delta f_n$ is the width of that bin, and $\Delta f_{tot}=\sum_n\Delta f_n$ is the total bandwidth. 
The lower and upper frequency limits are selected for the region of interest, taking into account \mi{f}{sw} and \mi{f}{res}; the equation above uses the band from \SI{10}{MHz} to \SI{150}{MHz}.

\subsubsection*{Amplitude of Resonance Frequency {\boldmath$\hat{V}_\mathrm{res}$}}
Peaks in the emission spectrum may originate from the source spectrum of \mi{v}{DS} or from resonances in the transfer function introduced by passive components and layout parasitics. 
The resonance FOM, $\hat{V}_\mathrm{res}$, quantifies the spectral component at a selected resonance frequency.
Most of the time, this will be the main resonance, which is created by the oscillation between the output capacitance of the turned-off device and the parasitic inductance of the commutation loop.

\subsubsection*{Summary of FOMs}
The considered FOMs and the information they provide are summarized in Tab. \ref{tab:FOMs}; their influence on the spectrum is illustrated in Fig. \ref{fig:FOM_in_spectrum}.


\begin{table}[ht]
	\centering
	\caption{Summary of all considered FOMs for the emitted-noise performance of a switching event.}
	\begin{tabular}{c|c|c}		
		FOM	& domain &  information \\
		\hline
		\mi{t}{10-90\%} & time & $2^\mathrm{nd}$ cut-off frequency \\
		$\sddt{v}_\mathrm{max}$ & time & correlates with \mi{V}{DC} \& \mi{t}{10-90\%}\\
		HGI & hybrid & difference to EMC optimum \\
		$\bar{V}_{\Delta f}$ & frequency & total noise emission \\
		$\hat{V}_\mathrm{res}$ & frequency & resonance \\
	\end{tabular}
	\label{tab:FOMs}
\end{table}

\begin{figure}
	\centering\import{./figures/3_Methods/FOM/}{Diss_FOM_in_Spectrum.pdf_tex}
	\caption{Influence of FOM in spectrum.}
	\label{fig:FOM_in_spectrum}
\end{figure}
\FloatBarrier


\subsubsection*{Time-Resolved Frequency-Domain Analysis}
When a parameter varies during a measurement, the frequency-domain FOMs can be evaluated over time to track the corresponding changes in the emission spectrum. 
A spectrum calculated over the entire record would combine measurements from different parameter values and obscure this evolution. 
Because the spectrum also depends on the switching duty cycle, it is not meaningful to calculate these FOMs from an isolated transient. 
Instead, a moving window is used to calculate the spectrum over $n$ switching periods while separating consecutive events, as illustrated in \fref{fig:Moving_FD_FOM}, and the resulting signal segment is repeated to provide sufficient length for the FFT. 
This procedure can be applied either to the original waveform or to event-separated signals generated by the snap-and-mirror method.

\begin{figure}
	\centering\import{./figures/3_Methods/FOM/}{Moving_FD_FOM.pdf_tex}
	\caption{Moving window for calculation of FD FOM over time with a window length of n = 1.}
	\label{fig:Moving_FD_FOM}
\end{figure}

